\subsection{Epilogue}

Professor Evjen is a voice lost to history. You will find no monument or statue. What you will find is living memory and a LFC footnote regarding the contents of this book that has become known as the ``Evjen Controversy.'' 

Evjen emerges at a particularly difficult time for the FLC. Sverdrup and Oftedal are dead. The cultural default of the Norwegian ethnicity is waning against a powerful American patriotic current fueled by the relentless advance of technology, westward expansion, and events such as WWI. 

We are all Americans, and we are Americans first. If that wasn't true in 1918 it accelerated in the crucible of WWII. The unfortunate result was the loss of a living Norwegian language in America. We nearly lost access to Sverdrup, Oftedal and Evjen in the process. We can thank the Sverdrup Society for carrying the torch. Recent efforts to digitize the world's libraries and the rise of AI-assisted research allow us to revisit the wisdom of the FLC founders.  

This change in temperature is reflected in the AFLC Rules for Work as adopted in 1967. Originally rule \#2 stated ``Its purpose is to labor for the vivification and liberation of all Norwegian Lutheran congregations...'' this was changed to ``Its aim shall be to work toward making Lutheran congregations free and living, so that, according to their calling and ability...''

I would argue that the ``Evjen Controversy'' is not so much about Evjen as the sudden, unexpected, and undisciplined transfer of power in the FLC. In the founders’ absence, people naturally question who they are and the meaning of their mission exemplified in the founding principles. If you will allow the imperfect analogy, this is Israel wandering in the desert without Moses. The promise was given but people lost direction. Many returned to the larger synodical body. Others soldiered forward with unquestioned allegiance to Sverdrup and Oftedal, forgetting that the founders are witnesses to a principle and not a substitute for the real thing. 

Evjen with his systemizing probe of at the founding principles was salt in the wound. His actions were unwelcome in 1918 and remain challenging today in 2026. In fact, it may be harder in 2026 as the polemic language seems overly harsh and foreign to our ears.   

Returning once again to Evjen, this is what we do know. Shortly after writing \textit{The Position} in Folkebladet, he was dismissed from the seminary. Tragically, I've found no additional Evjen writings in Folkebladet. His name emerges again as head of Mayville, North Dakota. There he wrote a piece titled \textit{The Teachers College: Its Place in the Educational System}.\footnote{John O. Evjen, The Teachers College: Its Place in the Educational System, State Normal School Quarterly Bulletin, new ser., no. 12 (December 1920): 17.} You may appreciate this quote where we see the same Evjen pattern all over again. He suffers no fools and drives straight at the heart of the problem. %He calls forth the snakes under the carpet.


\begin{quote}
A WEAKNESS IN HIGHER EDUCATION

The trouble with many of our colleges to-day is this: They have aspired to be universities, they have copied the English system, making the college the zenith institution of education. And yet, the copying has been so imperfect that the average American B. A. lacks three years of being equal to an Oxford B. A. For some inexplicable reason our colleges have built a German postgraduate department on top of the collegiate,
and called the entire confusion a ``university.''
\end{quote}

Ironically, this is another live-wire topic that has been exposed in today's America. In that sense, Evjen refuses to be silent and calls to us a century later.   

As Israel is forever called to wrestle with God, so too the AFLC must wrestle with the Principles beginning with P1, not as the slogan ``spirit and life'' but as a living vocational obligation that calls all hands to work for the Kingdom. Therefore, pray for wisdom and guidance as you walk your earthly vocation as an academic, seminarian, pastor, or layperson. We are all called to witness and live without apology under this confession: ``According to the Word of God, the congregation is the right form of the Kingdom of God on earth.''

\textit{Fail forward.}

